\chapter{ÁREA DE APLICACIÓN}
\label{cap:area-aplicacion}

\section{Antecedentes}
\label{sec:aa-antecedentes}

El área de aplicación del proyecto es el \emph{entrenamiento de fuerza y
acondicionamiento físico en gimnasios}. Un gimnasio comercial ofrece a sus
miembros acceso a un conjunto de equipamiento, generalmente distribuido en
zonas de máquinas de resistencia guiada, poleas, peso libre con bancos y
racks, y equipos cardiovasculares. La orientación que recibe el usuario
depende del gimnasio: en algunos casos hay un instructor de sala que atiende a
todos los presentes, y en la mayoría el usuario entrena solo y recurre a la
asesoría de un entrenador personal únicamente si puede pagarla.

En este contexto, el conocimiento necesario para entrenar de forma segura y
eficaz (cómo se ajusta y se usa cada máquina, qué músculos trabaja, cuántas
series y repeticiones realizar, cómo organizar la semana y cómo saber si se
está progresando) se transmite de manera informal: por imitación de otros
usuarios, por videos genéricos o por rutinas copiadas que no consideran el
nivel ni el objetivo de la persona. Esta situación es la que conduce a las
barreras descritas en la sección~\ref{sec:problema}.

\section{Descripción del área de aplicación}
\label{sec:aa-descripcion}

\subsection{Equipamiento de un gimnasio}

Para construir la guía de máquinas se tomó como referencia el catálogo de un
proveedor nacional de equipamiento para gimnasios. A partir de él se
identificaron 35 tipos de máquina representativos, agrupados en cuatro tipos
según su forma de uso (Tabla~\ref{tab:tipos-equipamiento}). Cada tipo agrupa
varios modelos comerciales, que en la aplicación se muestran como variantes
del mismo equipo.

\begin{table}[H]
\centering
\caption{Tipos de equipamiento considerados en el catálogo}
\label{tab:tipos-equipamiento}
\begin{tabularx}{\textwidth}{|L{3cm}|C{2.0cm}|Y|}
\hline
\rowcolor{encabezado} \textbf{Tipo} & \textbf{Cantidad} & \textbf{Descripción} \\
\hline
Máquina guiada & 23 & Equipo con trayectoria fija y placas de peso, por
ejemplo prensa de piernas, extensión de cuádriceps o máquina de pecho. \\ \hline
Banco y rack & 5 & Estructuras para entrenar con barra o mancuernas, como el
banco plano o el rack de sentadillas. \\ \hline
Cardio & 4 & Equipos de trabajo cardiovascular, como la cinta de correr o la
bicicleta estática. \\ \hline
Polea & 3 & Estaciones de cable con trayectoria libre, como el cruce de poleas
o el jalón al pecho. \\
\hline
\textbf{Total} & \textbf{35} & \\
\hline
\end{tabularx}
\fuente{Elaboración propia a partir del catálogo cargado en Cloud Firestore.}
\end{table}

\subsection{Actores}

En el área de aplicación intervienen los siguientes actores, que se
corresponden con los perfiles de usuario de la aplicación:

\begin{itemize}
    \item \textbf{Usuario principiante:} lleva menos de seis meses entrenando
    o es nuevo en el gimnasio. Necesita, sobre todo, aprender a usar las
    máquinas y seguir una rutina sencilla. Es el actor con mayor riesgo de
    abandono.
    \item \textbf{Usuario intermedio:} entrena con regularidad desde hace seis
    meses a dos años. Conoce las máquinas básicas y busca una rutina más
    estructurada y el seguimiento de su progreso.
    \item \textbf{Usuario avanzado:} entrena desde hace más de dos años. Arma
    sus propias rutinas y valora la planificación y los reportes.
    \item \textbf{Administrador del catálogo:} rol técnico, fuera de la
    aplicación, responsable de cargar y actualizar el contenido de las
    máquinas mediante scripts con credenciales de servicio.
\end{itemize}

Los objetivos de entrenamiento que el usuario puede declarar son cuatro:
\emph{perder peso}, \emph{ganar masa muscular}, \emph{mantenimiento} y
\emph{resistencia}.

\section{Procesos}
\label{sec:aa-procesos}

\subsection{Proceso actual}

La Figura~\ref{fig:proceso-actual} representa el proceso que sigue hoy un
usuario sin guía profesional, y la Tabla~\ref{tab:problemas-proceso} relaciona
cada paso con el problema que genera.

\begin{figure}[H]
\centering
\begin{tikzpicture}[node distance=0.45cm and 0.5cm,
  paso/.style={caja, minimum width=2.6cm, text width=2.5cm, font=\footnotesize},
  prob/.style={rectangle, draw=red!60!black, fill=red!5, dashed, rounded corners=2pt,
               text width=2.5cm, align=center, font=\scriptsize}]
  \node[paso] (a) {Se inscribe en el gimnasio};
  \node[paso, right=of a] (b) {Imita a otros o copia una rutina genérica};
  \node[paso, right=of b] (c) {Usa las máquinas sin ajuste ni técnica};
  \node[paso, right=of c] (d) {No registra su progreso};
  \node[prob, below=0.7cm of b] (pb) {Rutina sin relación con su objetivo ni nivel};
  \node[prob, below=0.7cm of c] (pc) {Riesgo de lesión y bajo resultado};
  \node[prob, below=0.7cm of d] (pd) {No percibe avances};
  \node[caja, fill=red!12, below=0.7cm of pc, minimum width=5cm] (z) {Desmotivación y abandono};
  \draw[flecha] (a)--(b); \draw[flecha] (b)--(c); \draw[flecha] (c)--(d);
  \draw[flechad] (b)--(pb); \draw[flechad] (c)--(pc); \draw[flechad] (d)--(pd);
  \draw[flecha] (pb) |- (z); \draw[flecha] (pc)--(z); \draw[flecha] (pd) |- (z);
\end{tikzpicture}
\caption{Proceso actual de entrenamiento sin guía profesional}
\label{fig:proceso-actual}
\end{figure}

\begin{table}[H]
\centering
\caption{Problemas del proceso actual y respuesta del sistema}
\label{tab:problemas-proceso}
\begin{tabularx}{\textwidth}{|Y|Y|C{1.9cm}|}
\hline
\rowcolor{encabezado} \textbf{Problema} & \textbf{Respuesta del sistema} & \textbf{Objetivo} \\
\hline
No sabe cómo se ajusta ni cómo se usa cada máquina. & Guía de máquinas con
instrucciones, errores comunes y prevención de lesiones; escáner con
cámara. & OE1 \\ \hline
No sabe ejecutar correctamente los ejercicios. & Animaciones de ejecución por
ejercicio, con músculo objetivo y nivel de dificultad. & OE2 \\ \hline
Su rutina no considera su objetivo ni su nivel. & Generador de rutinas basado
en reglas de prescripción del ejercicio. & OE3 \\ \hline
No organiza su semana de entrenamiento. & Plan semanal por días y
planificador de rutinas personalizadas. & OE4 \\ \hline
No controla su composición corporal. & Calculadora de IMC con historial de
mediciones. & OE5 \\ \hline
No registra su evolución. & Registro de peso, estatura e IMC en el tiempo. &
OE6 \\ \hline
No percibe sus avances. & Reporte gráfico de la tendencia del peso y
resumen de la variación. & OE7 \\
\hline
\end{tabularx}
\end{table}

\subsection{Proceso propuesto}

Con el sistema, el usuario sigue el proceso de la
Figura~\ref{fig:proceso-propuesto}. Después de registrarse, declara su
objetivo y su nivel; con esos datos obtiene una rutina semanal generada
automáticamente o arma la suya. Durante el entrenamiento consulta la guía de
cada máquina, ya sea buscándola en el catálogo o enfocándola con la cámara, y
periódicamente registra su peso para visualizar su progreso.

\begin{figure}[H]
\centering
\begin{tikzpicture}[node distance=0.55cm and 0.45cm,
  paso/.style={caja, minimum width=2.7cm, text width=2.6cm, font=\footnotesize}]
  \node[paso, fill=green!8] (a) {Registro o inicio de sesión};
  \node[paso, fill=green!8, right=of a] (b) {Onboarding: objetivo y nivel};
  \node[paso, fill=green!8, right=of b] (c) {Rutina generada o personalizada};
  \node[paso, fill=green!8, right=of c] (d) {Entrenamiento del día};
  \node[paso, below=1.1cm of d] (e) {Guía de la máquina (catálogo o escáner)};
  \node[paso, left=of e] (f) {Registro de peso e IMC};
  \node[paso, left=of f] (g) {Reporte de progreso};
  \node[paso, left=of g] (h) {Ajuste de la rutina};
  \draw[flecha] (a)--(b); \draw[flecha] (b)--(c); \draw[flecha] (c)--(d);
  \draw[flecha] (d)--(e); \draw[flecha] (e)--(f); \draw[flecha] (f)--(g);
  \draw[flecha] (g)--(h); \draw[flecha] (h) -- ++(0,0.9) -| ($(c.south)+(0,0)$);
\end{tikzpicture}
\caption{Proceso propuesto de entrenamiento con el sistema}
\label{fig:proceso-propuesto}
\end{figure}

\section{Definiciones}
\label{sec:aa-definiciones}

\begin{description}[style=nextline, leftmargin=0.6cm, font=\normalfont\bfseries]
  \item[Repetición] Ejecución completa de un movimiento, desde la posición
  inicial hasta la final y el retorno.
  \item[Serie] Conjunto de repeticiones consecutivas realizadas sin descanso.
  \item[Descanso] Tiempo de recuperación entre series, expresado en segundos.
  \item[Rutina] Plan de entrenamiento que define, para cada día de la semana,
  los ejercicios a realizar con sus series, repeticiones y descansos.
  \item[Plan semanal] Distribución de los días de entrenamiento de una rutina,
  cada uno con un enfoque muscular (por ejemplo, ``Pecho, Hombros y Brazos'').
  \item[Rutina activa] Rutina que el usuario sigue actualmente y que se muestra
  en la pantalla de inicio.
  \item[Grupo muscular principal y secundario] Músculos que realizan la mayor
  parte del trabajo en un ejercicio y músculos que colaboran en él,
  respectivamente.
  \item[Máquina guiada] Equipo de resistencia con trayectoria fija que limita
  el movimiento a un plano determinado.
  \item[Variante] Modelo comercial concreto de un tipo de máquina.
  \item[Ficha de máquina] Conjunto de información de una máquina en la
  aplicación: descripción, músculos, guía de uso, errores comunes, prevención
  de lesiones, variantes y ejercicios asociados.
  \item[Onboarding] Proceso inicial en el que el usuario recién registrado
  declara su objetivo y su nivel de experiencia.
  \item[IMC] Índice de masa corporal (Ecuación~\ref{eq:imc}).
  \item[Catálogo] Conjunto de datos compartido por todos los usuarios (máquinas
  y ejercicios), de solo lectura para la aplicación.
\end{description}
